Este trabalho propõe uma abordagem analítica sobre métodos de processamento 
implementados em sistemas baseados em tecnologias \gls{iot} com o objetivo 
de aplicação da técnica de \gls{retrofitting} em instrumentos de medição analógicos.
Tal temática, tem como alicerce a lacuna literária identificada por \citeonline{Hasan2019}, 
onde carecem trabalhos que apresentam informações úteis para a etapa de desenvolvimento arquitetural 
de tais ferramentas.

Diante de tal contexto, revela-se a importância acadêmica deste trabalho, 
apresentando de forma científica comparações entre os métodos de processamento, 
definindo limites e estipulando vantagens na adoção de determinado método. E ainda, 
este trabalho tem como pauta discernir sobre os temas: 
desenvolvimento de sistemas baseados em tecnologia \gls{lora},
aplicação de algoritmos de packet loss (Perdas de Pacote) e
uplink de dados segmentados em pacotes via \gls{lorawan}.

O desenvolvimento deste trabalho, tem por foco a disponibilização de insumos para 
tomada de decisões no desenvolvimento de sistemas \gls{iot} que empregam protocolos de 
rede de baixa largura de banda.
